% Pre-ambulo
\documentclass[a4paper, 12pt]{abnt}

\usepackage[brazil]{babel}
\usepackage[T1]{fontenc}
\usepackage[utf8]{inputenc}
\usepackage{dsfont}
\usepackage{amssymb,amsmath}
\usepackage{multirow}
\usepackage[alf]{abntcite}
\usepackage[pdftex]{color, graphicx}
\usepackage{colortbl}
\usepackage{url}
\usepackage{abnt-alf}
\usepackage{abntcite}
\usepackage{algorithm}
\usepackage{algorithmic}
%\usepackage{alg}
%\usepackage{hyperref}


% Redefinicao de instrucoes
\floatname{algorithm}{Algoritmo}
\renewcommand{\algorithmicrequire}{\textbf{Entrada:}}
\renewcommand{\algorithmicensure}{\textbf{Saída:}}
\renewcommand{\algorithmicend}{\textbf{fim}}
\renewcommand{\algorithmicif}{\textbf{se}}
\renewcommand{\algorithmicthen}{\textbf{então}}
\renewcommand{\algorithmicelse}{\textbf{senão}}
\renewcommand{\algorithmicfor}{\textbf{para}}
\renewcommand{\algorithmicforall}{\textbf{para todo}}
\renewcommand{\algorithmicdo}{\textbf{faça}}
\renewcommand{\algorithmicwhile}{\textbf{enquanto}}
\renewcommand{\algorithmicloop}{\textbf{loop}}
\renewcommand{\algorithmicrepeat}{\textbf{repetir}}
\renewcommand{\algorithmicuntil}{\textbf{até que}}
\renewcommand{\algorithmiccomment}[1]{\% #1}


% Definicao da lista de simbolos
% \simb[entrada na lista de simbolos]{simbolo}:
% Escreve o simbolo no texto e uma entrada na lista de simbolos.
% Se o parametro opcional e omitido, usa-se o parametro obrigatorio.
\newcommand{\simb}[2][]
{%
	\ifthenelse{\equal{#1}{}}
	{\addcontentsline{los}{simbolo}{#2}}
	{\addcontentsline{los}{simbolo}{#1}}#2
}
% Para aceitar comandos com @ (at) no nome
\makeatletter 
% \listadesimbolos: comando que imprime a lista de simbolos
\newcommand{\listadesimbolos}
{
	\pretextualchapter{Lista de símbolos}
	{\setlength{\parindent}{0cm}
	\@starttoc{los}}
}
% Como a entrada sera impressa
\newcommand\l@simbolo[2]{\par #1}
\makeatother

\newcommand{\ignore}[1]{}

% Definicao da lista de abreviaturas e siglas
% \abrv[entrada na lista de simbolos]{abreviatura}:
% Escreve a sigla/abreviatura no texto e uma entrada na lista de abreviaturas e siglas.
% Se o parametro opcional e omitido, usa-se o parametro obrigatorio.
\newcommand{\abrv}[2][]
{%
	\ifthenelse{\equal{#1}{}}
	{\addcontentsline{loab}{abreviatura}{#2}}
	{\addcontentsline{loab}{abreviatura}{#1}}#2
}
% Para aceitar comandos com @ (at) no nome
\makeatletter 
% \listadeabreviaturas: comando que imprime a lista de abreviaturas e siglas
% \newcommand{\listadeabreviaturas}
% {
% 	\pretextualchapter{Lista de abreviaturas e siglas}
% 	{\setlength{\parindent}{0cm}
% 	\@starttoc{loab}}
% }
% Como a entrada sera impressa
% \newcommand\l@abreviatura[2]{\par #1}
\makeatother


% \listofalgorithms: comando que imprime a lista de algoritmos
\renewcommand{\listalgorithmname}{Lista de algoritmos}


% Hifenização de palavras feita de forma incorreta pelo LaTeX
\hyphenation{PYTHON ou-tros}


% Inicio do documento
\begin{document}

	\frenchspacing
	
	% Capa (arquivo Includes/Capa.tex)
	% Capa
% Proteção externa do trabalho e sobre a qual se imprimem as informações indispensáveis 
% à sua identificação.

% Especificação da capa
\begin{titlepage}
	\begin{center}
		
		% Cabeçalho (não deve ser modificado)
		% Contém o brasão da Universidade, o logotipo do Departamento, além dos dados
		% relacionados à vinculação do aluno (Universidade, Centro, Departamento e Curso)
		\begin{minipage}{2cm}
			\begin{center}
				\includegraphics[width=1.7cm, height=2.0cm]{Imagens/Brasao-UFRN.jpg}
			\end{center}
		\end{minipage}
		\begin{minipage}{11cm}
			\begin{center}
				\begin{espacosimples}
					{\small \textsc{Universidade Federal do Rio Grande do Norte}			\\
							  \textsc{Centro de Ciências Exatas e da Terra}						\\
							  \textsc{Departamento de Informática e Matemática Aplicada}	\\
							  \textsc{Bacharelado em Engenharia de Software}}
				\end{espacosimples}
			\end{center}
		\end{minipage}
		\begin{minipage}{2cm}
			\begin{center}
				\includegraphics[width=1.8cm, height=1.5cm]{Imagens/Logotipo-DIMAp.jpg}
			\end{center}
		\end{minipage}
			
		\vspace{6cm}
						
		% Título do trabalho
		{\setlength{\baselineskip}%
		{1.3\baselineskip}
		{\LARGE \textbf{Ferramenta de Suporte à Catalogação de Implementações de Funcionlidades de Software}}\par}
			
		\vspace{4cm}
			
		% Nome do aluno (autor)
		{\large \textbf{Igor Marques da Silva}}
						
		\vspace{7cm}
		
		% Local da instituição onde o trabalho deve ser apresentado e ano de entrega do mesmo
		Natal-RN\\Dezembro 2014
	\end{center}
\end{titlepage}

	% Folha de rosto (arquivo Includes/FolhaRosto.tex)
	% Folha de rosto
% Contém os elementos essenciais à identificação do trabalho.

% Título, nome do aluno e respectivo orientador e filiação
\titulo{\Large{Título do trabalho}}
\autor{Igor Marques da Silva}
\orientador[Orientador(a)]{\par Nome e titulação do(a) professor(a) orientador(a)}
\instituicao
{
	Universidade Federal do Rio Grande do Norte -- UFRN \par 
	Departamento de Informática e Matemática Aplicada -- DIMAp
}
	
% Natureza do trabalho (não deve ser modificada)
\comentario
{
	Monografia de Graduação apresentada ao Departamento de Informática e Matemática Aplicada do 
	Centro de Ciências Exatas e da Terra da Universidade Federal do Rio Grande do Norte como
	requisito parcial para a obtenção do grau de bacharel em Ciência da Computação.
}
		
% Local e data
\local{Natal-RN}
\data{Dezembro de 2014}
	
\folhaderosto	
	
	% Folha de aprovacao (arquivo Includes/FolhaAprovacao.tex)
	\include{Includes/FolhaAprovacao}	
	
	% Dedicatoria (arquivo Includes/Dedicatoria.tex)
	\include{Includes/Dedicatoria}
	
	% Agradecimentos (arquivo Includes/Agradecimentos.tex)
	\include{Includes/Agradecimentos}
   
   % Epigrafe (arquivo Includes/Epigrafe.tex)
	\include{Includes/Epigrafe}
	
	% Resumo em língua vernacula (arquivo Includes/Resumo.tex)
	\include{Includes/Resumo}
	
	% Abstract, resumo em língua estrangeira (arquivo Include/Abstract.tex)
	\include{Includes/Abstract}
	
	% Lista de figuras
% 	\listoffigures

	% Lista de tabelas
% 	\listoftables
	
	% Lista de abreviaturas e siglas
% 	\listadeabreviaturas
	
	% Lista de símbolos
% 	\listadesimboenlos
	
	% Lista de algoritmos (se houver)
	% Devem ser incluídos os pacotes algorithm e algorithmic
	% \listofalgorithms
	
	% Sumário
	\sumario

	% Parte central do trabalho, englobando os capítulos que constituem o mesmo
	% Os referidos capítulos devem ser organizados dentro do diretório "Capítulos"

	\chapter{Introdução}

No contexto de trabalhadores do conhecimento, como na área de desenvolvimento de software, é diária a troca de conhecimento com o objetivo de se obter um melhor desempenho para a organização inteira~\cite{Druker1993, Wiig2003}. A gerência de conhecimento lida com o reuso de conhecimento em  suas diferentes formas, como: design de código, requisitos, modelos, dados, padrões e lições aprendidas~\cite{Levy2009}. Um dos casos de reuso é o conhecimento de como se deu a implementação de uma determinada funcionalidade em um determinado projeto de software. Exemplo: diferentes projetos podem implementar cadastro de usuário via Google ou Facebook, cada um com certas especificidades.

É comum também que essa implementação em diferentes contextos seja feita através de abordagens \textit{ad-hoc} (reimplementação completa de trechos de código com poucas modificações e sem modularizacão visando reuso)~\cite{SangMok2011}. Nesse caso, desenvolvedores mais experientes em um determinado projeto ou que já implementou tal funcionalidade tendem a atuar como mentores~\cite{CubraniC2004} e tal ato, como atividade de gerência de conhecimento, acarreta no despendimento de recursos, principalmente dos mentores~\cite{Wiig2003}. 

Atualmente, as maneiras mais comuns de troca de informações entre desenvolvedores (incluindo, obviamente, a mentoria) são via oral, escrita ou repasse de referências (documentação, links externos, etc)~\cite{Storey2014, Olson2000, CubraniC2004}. Em alguns casos, pela própria natureza do meio, comunicação não se mantém registrada física ou virtualmente. Um dos exemplos de comunicação sem registro é a comunicação exclusivamente oral em que, ao fim do diálogo, as informações trocadas ficam apenas na memória dos envolvidos, sem maneiras de terceiros consultarem posteriormente aquelas informações~\cite{Olson2000}. 

Além disso, existe o fator da rotatividade de membros em equipes, sempre associado a custos de transferência de conhecimento e treinamento~\cite{Hall2008}, possuindo relação com sucesso ou fracasso de projetos de software~\cite{Hall2008}. 

Assim, a elaboração de uma ferramenta capaz de agregar referências de código e tarefas a soluções pode trazer enormes benefícios a equipes de desenvolvimento~\cite{CubraniC2004}. A ferramenta atuaria como um catálogo, agregando referências de código, informações externas e comentários fornecidos por desenvolvedores da equipe com o intuito de auxiliar outros desenvolvedores a buscar em fontes da própria equipe como se deram implementações de funcionalidades semelhantes às que já foram feitas em outros projetos.

Este estudo propõe a elaboração de tal ferramenta com a participação de uma equipe real de desenvolvimento de software. A equipe em questão é a 4Soft\footnote{\url{http://www.4softjr.com.br}}, empresa júnior\footnote{\url{http://en.wikipedia.org/wiki/Junior_enterprise}} dos cursos de Engenharia de Software e Tecnologia da Informação da Universidade Federal do Rio Grande do Norte (UFRN\footnote{\url{http://www.ufrn.br}}). A empresa atua na área de desenvolvimento de software Web para clientes de diversos ramos e é formada exclusivamente por alunos dos cursos de Bacharelado em Engenharia de Software\footnote{\url{http://www.dimap.ufrn.br/pt/graduacao/engenharia-de-software/apresentacao}} e Bacharelado em Tecnologia da Informação\footnote{\url{http://www.imd.ufrn.br/curso_bacharelado.php}} da UFRN.

O estudo também prevê uma análise crítica dos impactos do uso da mesma no contexto de empresa. Dentre os impactos previstos estão a redução da necessidade de mentoria ou consulta para os casos de  reuso de funcionalidades em diferentes contextos. Tais impactos visam contribuir para a redução de gastos de recursos relacionados a tempo de implementação de funcionalidades, esforço de recuperação de informações, orientação de colaboradores e trabalho dos desenvolvedores que atuam como mentores dentro da equipe. 

O Capítulo 2 descreve em detalhes o objetivo pretendido pelo trabalho e a quais perguntas pretende-se responder. Em seguida, o Capítulo 3 aborda a metodologia utilizada para a execução do estudo e que atividades serão desempenhadas pelos envolvidos. O Capítulo 4, por fim, descreve o cronograma das atividades descritas no capítulo anterior. 
	
	\chapter{Objetivos}

% gerais e específicos
Este trabalho tem por finalidade analisar que influência uma ferramenta de documentação dinâmica de funcionalidades de software traz para equipes de desenvolvimento. Além disso, objetiva verificar como tal ferramenta pode contribuir para a redução significativa de tempo despendido para a explicação repetitiva de tarefas já executadas anteriormente e como a troca de conhecimento entre membros pode fluir de maneira melhor e com baixo custo.

%falar do o que é uma documentacao dinamica


%falar quais os requisitos e os beneficios 

% perguntas de pesquisa

Assim, este trabalho visa responder às seguintes perguntas de pesquisa:

\begin{enumerate}
\item Como tal ferramenta pode contribuir para a documentação dinâmica de funcionalidades?
\item Como estimular o uso dessa ferramenta em equipes de desenvolvimento?
\item Que benefícios tal ferramenta provê a equipes de desenvolvimento?
\item Quais limitações tal ferramenta possui?

\end{enumerate}


   
	
	\chapter{Métodos}

% O que será feito, como será feito, … para abordar o problema

Este trabalho prevê o desenvolvimento da ferramenta de documentação de implementação funcionalidades. Nela, será possível vincular recursos do repositório do projeto no GitHub (issues, commits, etc) à requisitos e suas implementações, bem como outros referenciais (links para perguntas no Stack Overflow, desenhos, por exemplo) de forma a gerar um guia ou tutorial de como realizar tal implementação novamente no futuro.

A empresa 4Soft terá participação significativa em todo o estudo, desde a concepção até no uso da ferramenta.

% metodo: entrevista, questionário de satisfacao. Focar em pensar em como responder as perguntas

Os procedimentos para o estudo são os seguintes:

\section {Estudo de aplicações existentes}

Será feita uma busca por aplicações semelhantes. Suas limitações serão analisadas de modo que a ferramenta proposta as supra.

\section {Inquérito contextual}

Entrevistas e seções de brainstorm serão feitas com os participantes da empresa júnior mencionada. Será analisado como se dá seu processo de trabalho e como pode se dar o fluxo de atividades na ferramenta através de entrevistas e aplicação de questionários de satisfação.

\section {Sessões de interpretação da equipe}

Reuniões com a equipe de pesquisadores da UFRN que trabalharão no projeto serão feitas para definir o escopo da ferramenta, bem como seu se dará seu design e implementação

\section {Prototipação e implementação da ferramenta}

Nesta etapa, inicialmente, protótipos de baixa fidelidade serão elaborados. Posteriormente, serão expostos a todos os participantes e seu feedback será colhido e analisado. A partir daí, protótipos de maior fidelidade serão elaborados de maneira iterativa e incremental.

\section {Implantação e observação do uso da ferramenta}

A ferramenta então estará disponível para uso nos contextos descritos na seção anterior deste trabalho. A adesão dos desenvolvedores a ferramenta terá analisada nesta etapa, bem como seu uso (quantidade de artefatos de documentação criados, por exemplo) monitorado.

%resposta das peguntas 1 e 2


\section {Avaliação da ferramenta}

Ao fim do período anterior, uma nova bateria de entrevistas e aplicação de questionários de satisfação serão realizadas para avaliar qualitativamente como se deu a utilização da ferramenta, bem como se deram os efeitos de seu uso. 

%resposta das peguntas 3 e 4

A versão final do trabalho será então redigida.

	
	\chapter{Plano de Trabalho}

% Tarefas ou Atividades (bem macro)
% Cronograma

O cronograma de atividades se dará conforme o seguinte planejamento:

\begin{enumerate}

\item Estudo de aplicações existentes: Dezembro/2014

\item Inquérito contextual: Janeiro/2015

\item Sessões de interpretação da equipe: Janeiro/2015

\item Prototipação e implementação da ferramenta: Janeiro/2015 a Maio/2015

\item Implantação e observação do uso da ferramenta: Junho/2015 a Setembro/2015

\item Avaliação da ferramenta: Outubro/2015 a Novembro/2015

\end{enumerate}

		

	% Bibliografia (arquivo Capitulos/Referencias.bib)
	\bibliography{Capitulos/Referencias}
	\bibliographystyle{abnt-alf}
	
	% Apêndice A (arquivo Includes/ApendiceA)
	\include{Capitulos/ApendiceA}
	
	% Anexo A (arquivo Includes/AnexoA)
	\include{Capitulos/AnexoA}
	
	% Página em branco
	\newpage

\end{document}